% Every clinical prognostic index against OPERA at both ends. Rows are written by
% opera.run.build_paper_numbers (clinical_index_rows.tex) from the aggregate
% clinical-scores outputs; the placeholder row is shown until they exist.
\begin{table}[p]
\centering
\caption{\textbf{\opera{} against every clinical prognostic index at both ends.}
One row per index and the population it applies to (as in
\cref{fig:clinical-index-summary}), on test patients with the index recorded at
diagnosis, for the endpoint the index was developed for at the stated horizon.
Every difference is oriented so that positive values favour \opera{}.
\textbf{Equal group sizes}: \opera{}'s risk groups take the sizes of the index's
categories in the same patients; (a) survival of \opera{}'s lowest group minus
the index's lowest category, and (b) survival of the index's highest category
minus \opera{}'s highest group, in percentage points, with 95\% paired bootstrap
intervals that repeat the size matching in each replicate. \textbf{Equal risk}:
cutoffs fixed on training data; (c) additional patients in \opera{}'s low-risk
group, selected as the largest group with survival no lower than the index's
lowest category, as a percentage of patients and as a count ($n$); (d) additional
events captured per 1,000 patients by \opera{}'s high-risk group at the index's
specificity (inverse probability of censoring weighting). Intervals are
conditional on the frozen cutoffs. Bold: interval above zero; in colour:
interval below zero. n.e., not estimable; reasons are given where a whole
comparison could not be made. Counts across indices:
\result{clinical_low_size_atleast}{k of n} indices at least as good at the low
end and \result{clinical_high_size_atleast}{k of n} at the high end (estimate
$\geq0$; analysis-lock rule), \result{clinical_low_size_better}{k of n} and
\result{clinical_high_size_better}{k of n} better (interval above zero).}
\label{tab:clinical-index-summary}
\scriptsize
\setlength{\tabcolsep}{3.5pt}
\renewcommand{\arraystretch}{1.08}
\resizebox{\textwidth}{!}{%
\begin{tabular}{@{}lllrrrrrr@{}}
\toprule
& & & & \multicolumn{2}{c}{Equal group sizes} & \multicolumn{3}{c}{Equal risk} \\
\cmidrule(lr){5-6}\cmidrule(l){7-9}
Disease & Index & Endpoint & Patients & (a) Lowest group, pp & (b) Highest group, pp
  & (c) Low risk, \% & $n$ & (d) Events per 1,000 \\
\midrule
\IfFileExists{\resultnumbersdir clinical_index_rows.tex}{\input{\resultnumbersdir clinical_index_rows.tex}}{%
\placeholder{Disease} & \placeholder{Index} & \placeholder{Endpoint} & \placeholder{N}
  & \placeholder{x} & \placeholder{x} & \placeholder{x} & \placeholder{n} & \placeholder{x} \\
\bottomrule}
\end{tabular}}
\end{table}
